% Volume VIII: Computational Neuroscience Mathematics
% Continuity note for Chapter 49.
% Previous dependencies: Chapter 6's eigenvalues and quadratic forms; Chapter 8's ODE preview; Chapter 18's entropy, KL divergence, Fisher information, and information geometry; Chapter 22's variational principles; Chapter 25's fixed points and stability; Chapter 29's networks when available; Chapter 48's population state spaces and latent trajectories.
% New durable concepts: recurrent rate equations; fixed points; linearized stability; attractors and basins; Hopfield energy; Ising-style binary units; Gibbs and Boltzmann distributions; partition functions; free energy as energy-entropy tradeoff; mean-field self-consistency; oscillators, phase, synchrony, Kuramoto order parameter; neural field equations; static versus dynamic mean field.
% Forward hooks: Chapter 50 uses attractor and energy ideas cautiously in predictive coding and active inference; Chapter 51 uses circuit models as mechanistic hypotheses; Chapter 52 compares biological and artificial recurrent networks; Chapter 54 summarizes statistical-physics and dynamical-system identities.
% Notation commitments: Recurrent activity is x(t) or r(t) in \mathbb R^N; connectivity is W; nonlinearities are \phi; binary states are s_i in \{-1,+1\}; energy is E(s); inverse temperature is \beta; partition function is Z; neural field activity is u(x,t) with spatial coordinate x.
\section{Chapter 49: Neural Circuits, Attractors, Oscillations, and Neural Fields}
\label{ch:neural-circuits-attractors-oscillations-neural-fields}

Chapters 47 and 48 treated neural population activity as data: responses encode variables, and activity vectors trace geometry. Chapter 49 turns to mechanisms. A recurrent circuit is not just a source of responses; it is a dynamical system whose feedback can stabilize memories, amplify inputs, generate oscillations, and support spatial activity patterns.

The central mathematical problem is to move from many interacting units to understandable collective behavior. Dynamical-systems tools describe fixed points, stability, and oscillations. Statistical-physics tools describe energy landscapes, Gibbs distributions, partition functions, free energy, and mean-field approximations. Neural-field tools pass from discrete populations to continuous space. None of these tools proves a biological story by itself, but each gives a precise model class for circuit hypotheses.

\paragraph{Chapter problem.}
How do recurrent interactions create stable states, attractor memories, rhythms, spatial fields, and mean-field descriptions in neural and artificial networks?

\paragraph{Section map.}
Section 49.1 introduces recurrent rate circuits and linear stability. Section 49.2 develops attractor networks, Hopfield energies, Gibbs distributions, partition functions, and free energy. Section 49.3 defines oscillations, phase coupling, and synchrony. Section 49.4 introduces neural fields and mean-field approximations, including static versus dynamic mean field.

\subsection{49.1 Recurrent neural circuits}

\paragraph{Motivating question.}
How does feedback connectivity determine whether population activity decays, amplifies, stabilizes, or switches state?

\paragraph{Beginner setup.}
A recurrent circuit feeds activity back into itself. If \(x_i(t)\) is the activity of unit \(i\), then future activity depends on current activity of many units. The simplest continuous-time rate model is
\[
\tau \dot x=-x+W\phi(x)+u(t).
\]
The leak term \(-x\) pulls activity down. The recurrent term \(W\phi(x)\) sends nonlinear activity through the connectivity matrix. The input \(u(t)\) drives the circuit from outside.

Recurrent circuits solve a memory and computation problem. Feedforward responses disappear when the input disappears. Recurrent feedback can maintain a state, integrate evidence, select among alternatives, or generate patterned dynamics.

\paragraph{Formal setup.}
Let
\[
x(t)\in\mathbb R^N
\]
be the vector of neural activities, \(W\in\mathbb R^{N\times N}\) a recurrent weight matrix, \(u(t)\in\mathbb R^N\) an external input, and
\[
\phi:\mathbb R\to\mathbb R
\]
a scalar nonlinearity applied elementwise. A standard rate equation is
\[
\boxed{
\tau \frac{dx}{dt}=-x+W\phi(x)+u(t).
}
\]
Here \(\tau>0\) is a time constant. A fixed point for constant input \(u\) is a vector \(x^\star\) satisfying
\[
\boxed{
x^\star=W\phi(x^\star)+u.
}
\]
To study local stability, write \(x=x^\star+\delta x\). If \(D_\phi(x^\star)\) is the diagonal matrix with entries \(\phi'(x_i^\star)\), then the linearized dynamics are
\[
\tau\frac{d\,\delta x}{dt}
=\left[-I+W D_\phi(x^\star)\right]\delta x.
\]
The fixed point is locally stable if all eigenvalues of
\[
J=\frac{1}{\tau}\left[-I+W D_\phi(x^\star)\right]
\]
have negative real parts.

\paragraph{Core intuition.}
Feedback is useful because it can preserve and transform activity. It is dangerous because it can also amplify noise or create instability. Linearization gives the first local diagnostic. If a perturbation is pushed back toward the fixed point, the state is stable. If a perturbation grows, the fixed point is unstable.

The eigenvectors of the effective recurrent matrix give directions of perturbation. The eigenvalues tell whether those perturbations shrink, grow, or rotate. This is the same mathematical habit as Chapter 25's stability analysis, now applied to a neural circuit.

\paragraph{Main mechanism: scalar recurrent gain.}
For one unit,
\[
\tau \dot x=-x+w\phi(x)+u.
\]
A fixed point satisfies
\[
x^\star=w\phi(x^\star)+u.
\]
Linearizing around \(x^\star\) gives
\[
\tau \dot{\delta x}
=\left[-1+w\phi'(x^\star)\right]\delta x.
\]
Thus the fixed point is locally stable when
\[
\boxed{
w\phi'(x^\star)<1.
}
\]
The product \(w\phi'(x^\star)\) is the local recurrent gain. If the loop gain exceeds the leak, small deviations grow.

\paragraph{Worked examples.}
\emph{Small numerical example.}
Let \(\tau=10\) ms, \(w=1.5\), and suppose \(\phi'(x^\star)=0.4\). Then
\[
w\phi'(x^\star)=0.6<1,
\]
so the fixed point is locally stable. If instead \(w=3\) with the same slope, then \(w\phi'(x^\star)=1.2>1\), and the scalar fixed point is locally unstable.

\emph{Computational neuroscience example.}
An excitatory-inhibitory cortical microcircuit can be summarized by two variables \(x_E,x_I\), with weights describing E-to-E, E-to-I, I-to-E, and I-to-I feedback. Stability depends on the balance between recurrent excitation and inhibitory feedback. Too much effective excitation can create runaway activity; too much inhibition can silence or oscillate the circuit depending on delays and gains.

\emph{AI example.}
Recurrent neural networks and residual networks also contain feedback-like iterative computations. The spectral radius or effective Jacobian of the recurrent update influences whether gradients and activations explode, vanish, or remain stable over time.

\paragraph{Common misconceptions and failure modes.}
A connectivity matrix alone does not determine dynamics without the nonlinearity, input, and operating point. Stability is local unless a global argument is supplied. A stable fixed point is not automatically desirable; it may be an uninformative silent state. A biological circuit can have delays, spiking noise, adaptation, and plasticity not captured by a simple rate equation.

\paragraph{Exercises.}
\begin{enumerate}[label=\textbf{49.1.\arabic*.}, leftmargin=*]
\item \textbf{Mechanical.} For \(\tau\dot x=-x+w\phi(x)+u\), write the fixed-point equation and the local stability condition in terms of \(w\phi'(x^\star)\).
\item \textbf{Proof or derivation.} Derive the linearized matrix \(\tau\dot{\delta x}=(-I+WD_\phi(x^\star))\delta x\) for the \(N\)-unit rate model.
\item \textbf{Numerical/computational.} If \(W=\begin{bmatrix}1&-2\\1&-1\end{bmatrix}\), \(D_\phi=0.5I\), and \(\tau=1\), compute the effective linearization matrix \(J=-I+WD_\phi\).
\item \textbf{Interpretation.} In an excitatory-inhibitory circuit, why can inhibitory feedback stabilize recurrent excitation?
\item \textbf{Debug the misconception.} A student says, "Positive recurrent weights always imply instability." Give a condition under which positive recurrent feedback remains locally stable.
\end{enumerate}

\subsection{49.2 Attractor networks}

\paragraph{Motivating question.}
How can a recurrent network store patterns so that noisy or partial inputs relax back to stable states?

\paragraph{Beginner setup.}
An attractor is a stable set of states toward which nearby trajectories move. In a memory model, an attractor can represent a stored pattern. If the current activity is a corrupted version of the pattern, the dynamics may return it to the clean memory. The set of initial states that flow to the attractor is its basin of attraction.

Attractor networks make the geometric idea of an energy landscape concrete. Activity moves downhill in energy until it reaches a local minimum. Hopfield networks are the classic example. They use binary units and symmetric weights so that asynchronous updates decrease an energy function.

\paragraph{Formal setup.}
Let binary units satisfy
\[
s_i\in\{-1,+1\},\qquad s=(s_1,\ldots,s_N).
\]
Assume symmetric weights \(W_{ij}=W_{ji}\) and no self-couplings \(W_{ii}=0\). With biases \(b_i\), define the Hopfield energy
\[
\boxed{
E(s)=-\frac12\sum_{i,j}W_{ij}s_i s_j-\sum_i b_i s_i.
}
\]
An asynchronous update chooses one unit \(i\) and sets
\[
s_i\leftarrow \operatorname{sign}\left(\sum_j W_{ij}s_j+b_i\right).
\]
If ties are handled consistently, this update never increases \(E(s)\).

The statistical-physics version assigns a Gibbs or Boltzmann distribution
\[
\boxed{
p_\beta(s)=\frac{1}{Z_\beta}\exp[-\beta E(s)],
\qquad
Z_\beta=\sum_{s\in\{-1,+1\}^N}\exp[-\beta E(s)].
}
\]
The constant \(Z_\beta\) is the partition function. It normalizes the distribution. Computing it exactly requires summing over \(2^N\) states, which is infeasible for large \(N\).

\paragraph{Core intuition.}
Energy is a scalar score for network states. Low-energy states are stable memories or likely configurations. Temperature controls randomness: high inverse temperature \(\beta\) concentrates probability near energy minima, while low \(\beta\) spreads probability more broadly.

The partition function is hard because it counts the contribution of every possible network state. This is the same normalization problem that appears in energy-based machine learning models and in Bayesian evidence calculations.

\paragraph{Main theorem: asynchronous Hopfield updates decrease energy.}
Let only unit \(i\) change from \(s_i\) to \(s_i'\), with all other units fixed. Define the local field
\[
h_i=\sum_j W_{ij}s_j+b_i.
\]
The part of the energy depending on \(s_i\) is
\[
-s_i h_i
\]
because the symmetry and factor \(1/2\) count pairwise interactions once. Therefore the energy change is
\[
\Delta E=E(s_i')-E(s_i)=-(s_i'-s_i)h_i.
\]
The update \(s_i'=\operatorname{sign}(h_i)\) chooses the sign that makes \(-s_i'h_i\) as small as possible. Hence
\[
\Delta E\leq0.
\]
The network cannot cycle through states with strictly decreasing energy forever because there are finitely many states. It must eventually reach a state where every unit agrees with its local field, a local energy minimum.

\paragraph{Worked examples.}
\emph{Small numerical example.}
Consider two units with \(W_{12}=W_{21}=1\), \(b_1=b_2=0\). The energy is
\[
E(s)=-s_1s_2.
\]
The states \((1,1)\) and \((-1,-1)\) have energy \(-1\). The states \((1,-1)\) and \((-1,1)\) have energy \(+1\). Thus aligned states are attractors and anti-aligned states are higher-energy configurations.

\emph{Computational neuroscience example.}
Attractor models have been used for working memory, decision states, head direction, and pattern completion. In a working-memory interpretation, a bump or fixed point persists after the stimulus is removed. The biological claim requires evidence beyond the model, such as perturbation recovery, stability, and circuit mechanisms.

\emph{AI example.}
Boltzmann machines and modern energy-based models assign low energy to plausible configurations and high energy to implausible ones. Training often requires estimating gradients that involve model expectations under \(p_\beta(s)\), which is hard because of the partition function. Contrastive and variational methods approximate this normalization problem.

\paragraph{Free energy and the energy-entropy tradeoff.}
For any distribution \(q(s)\), define the variational free energy at temperature \(T=1/\beta\) as
\[
\boxed{
F(q)=\mathbb E_q[E(S)]-T H(q),
}
\]
where \(H(q)=-\sum_s q(s)\log q(s)\) in natural units. The first term prefers low-energy states. The entropy term prefers spreading probability across states. The Gibbs distribution minimizes this free energy:
\[
q^\star(s)=\frac{\exp[-E(s)/T]}{Z}.
\]
This is the precise meaning of free energy as an energy-entropy tradeoff. In variational inference and energy-based AI, one often restricts \(q\) to a tractable family and minimizes \(F(q)\) or an equivalent KL objective.

\paragraph{Common misconceptions and failure modes.}
An energy function is not physical energy unless the model specifies a physical system. It is a Lyapunov or probability-scoring function. A Hopfield energy decrease requires symmetry and asynchronous updates; arbitrary recurrent neural networks need not have an energy. A partition function being hard is not a philosophical issue; it is a concrete exponential sum or high-dimensional integral. Attractor language should not be used for any persistent activity unless stability and basin structure are specified.

\paragraph{Exercises.}
\begin{enumerate}[label=\textbf{49.2.\arabic*.}, leftmargin=*]
\item \textbf{Mechanical.} For the two-unit example with \(E(s)=-s_1s_2\), compute \(E\) for all four states.
\item \textbf{Proof or derivation.} Derive \(\Delta E=-(s_i'-s_i)h_i\) for one asynchronous Hopfield update.
\item \textbf{Numerical/computational.} For the two-unit example, compute the partition function \(Z_\beta\) at \(\beta=1\).
\item \textbf{Interpretation.} Explain why a corrupted memory pattern can return to a stored pattern in an attractor network.
\item \textbf{Debug the misconception.} A student says, "Every recurrent network has an energy function." State the assumptions that make the Hopfield energy argument work and why they are not automatic.
\end{enumerate}

\subsection{49.3 Oscillations and synchrony}

\paragraph{Motivating question.}
How can recurrent systems produce rhythms, and how can many rhythmic units become synchronized?

\paragraph{Beginner setup.}
An oscillation is repeated motion through state space. A scalar sinusoid
\[
x(t)=A\cos(\omega t+\varphi)
\]
has amplitude \(A\), angular frequency \(\omega\), and phase \(\omega t+\varphi\). Neural oscillations may appear as rhythmic membrane potentials, local field potentials, population firing rates, or latent trajectories.

Synchrony asks whether multiple oscillatory units keep a consistent phase relationship. Perfect synchrony means they rise and fall together. Phase locking means their phase difference stays approximately constant, even if it is not zero.

\paragraph{Formal setup.}
A phase oscillator is described by an angle
\[
\theta_i(t)\in[0,2\pi)
\]
with intrinsic frequency \(\omega_i\). A standard coupled phase model is the Kuramoto model:
\[
\boxed{
\dot\theta_i=\omega_i+\frac{K}{N}\sum_{j=1}^N\sin(\theta_j-\theta_i).
}
\]
The coupling strength is \(K\). The synchrony order parameter is
\[
\boxed{
Re^{i\psi}=\frac1N\sum_{j=1}^N e^{i\theta_j}.
}
\]
Here \(0\leq R\leq1\). If phases are spread uniformly around the circle, \(R\approx0\). If all phases are equal, \(R=1\).

\paragraph{Core intuition.}
Oscillations often arise when feedback does not simply return a perturbation directly to equilibrium, but rotates it through state space. Excitation, inhibition, delays, adaptation, and resonance can all produce rhythmic dynamics in different models.

Synchrony is a collective phenomenon. Weakly coupled oscillators may drift relative to each other. Strong enough coupling can overcome frequency differences and lock phases. The order parameter summarizes the group phase coherence.

\paragraph{Main mechanism: two-oscillator phase locking.}
For two Kuramoto oscillators,
\[
\dot\theta_1=\omega_1+\frac K2\sin(\theta_2-\theta_1),
\]
\[
\dot\theta_2=\omega_2+\frac K2\sin(\theta_1-\theta_2).
\]
Let \(\delta=\theta_2-\theta_1\). Then
\[
\dot\delta
=\dot\theta_2-\dot\theta_1
=\omega_2-\omega_1-K\sin\delta.
\]
A phase-locked solution satisfies \(\dot\delta=0\), so
\[
\boxed{
\sin\delta^\star=\frac{\omega_2-\omega_1}{K}.
}
\]
Such a solution exists only if
\[
|\omega_2-\omega_1|\leq K.
\]
Thus coupling must be strong enough relative to frequency mismatch.

\paragraph{Worked examples.}
\emph{Small numerical example.}
If \(\omega_2-\omega_1=1\) and \(K=2\), then
\[
\sin\delta^\star=1/2,
\]
so one locked phase difference is \(\delta^\star=\pi/6\). If \(K=0.5\), then \(|1|>0.5\), so no fixed phase difference exists in this simple model.

\emph{Computational neuroscience example.}
Gamma rhythms, theta rhythms, and beta rhythms are often studied through oscillatory population models and spectral analyses. A phase-locking analysis might ask whether spikes from one region occur at a consistent phase of a local field potential in another region. The mathematical question is phase relationship; the biological interpretation depends on task, measurement, and perturbation evidence.

\emph{AI example.}
Oscillatory and coupled-dynamical models appear in synchronization-based computation, recurrent networks, and some neuromorphic systems. More broadly, training dynamics in artificial networks can also show oscillations when optimization updates overshoot or interact with curvature, although those oscillations are in parameter space rather than neural phase.

\paragraph{Common misconceptions and failure modes.}
Not every peak in a spectrum implies a stable oscillator; finite data, filtering, and transients can create apparent rhythms. Synchrony does not automatically imply communication or causality. Phase is meaningful only after specifying the signal and extraction method. Averaging phases across trials with variable timing can hide or create apparent synchrony.

\paragraph{Exercises.}
\begin{enumerate}[label=\textbf{49.3.\arabic*.}, leftmargin=*]
\item \textbf{Mechanical.} For \(x(t)=3\cos(4t+\pi/2)\), identify amplitude, angular frequency, and phase offset.
\item \textbf{Proof or derivation.} Derive the two-oscillator phase-difference equation \(\dot\delta=\omega_2-\omega_1-K\sin\delta\).
\item \textbf{Numerical/computational.} Compute the order parameter \(R\) for phases \(0,0,0\), and then for phases \(0,2\pi/3,4\pi/3\).
\item \textbf{Interpretation.} What does phase locking mean for two neural populations with different intrinsic frequencies?
\item \textbf{Debug the misconception.} A student says, "High synchrony proves one population drives the other." Correct the statement using the distinction between phase relation and causality.
\end{enumerate}

\subsection{49.4 Neural field models}

\paragraph{Motivating question.}
How can many spatially arranged neurons be approximated by a continuous activity field, and how does mean-field theory replace interactions by self-consistent averages?

\paragraph{Beginner setup.}
A neural field model describes activity as a function of space and time:
\[
u(x,t).
\]
Here \(x\) is a position along cortex, a sensory map, or an abstract feature axis, and \(u(x,t)\) is population activity near that position. Instead of a finite matrix \(W_{ij}\), a neural field uses a connectivity kernel \(w(x-x')\) that describes how activity at \(x'\) influences activity at \(x\).

Mean-field theory has a related goal: replace many microscopic interactions by a smaller number of average quantities that must be self-consistent. In a static mean-field approximation, one solves for stationary average activity. In a dynamic mean-field approximation, one tracks time-dependent statistics such as autocorrelation or effective noise.

\paragraph{Formal setup.}
A standard one-dimensional Amari-style neural field equation is
\[
\boxed{
\tau\frac{\partial u(x,t)}{\partial t}
=-u(x,t)
+\int_{\mathcal X} w(x-x')\phi(u(x',t))\,dx'
+I(x,t).
}
\]
The kernel \(w\) describes spatial coupling, \(\phi\) is a firing-rate nonlinearity, and \(I\) is external input. A stationary field \(u^\star(x)\) satisfies
\[
u^\star(x)
=\int_{\mathcal X} w(x-x')\phi(u^\star(x'))\,dx'
+I(x).
\]

For a binary Ising-style population with states \(S_i\in\{-1,+1\}\), the mean-field approximation replaces neighbors \(S_j\) by their means
\[
m_j=\mathbb E[S_j].
\]
The self-consistency equation becomes
\[
\boxed{
m_i=\tanh\left(\beta\left[h_i+\sum_j J_{ij}m_j\right]\right).
}
\]
This is static mean field: the average field produced by the means must reproduce the same means.

\paragraph{Core intuition.}
Neural fields trade neuron labels for continuous position. A local-excitation, lateral-inhibition kernel can support bumps: localized regions of high activity surrounded by low activity. Such bumps are models for working memory, spatial attention, head direction, and feature maps.

Mean-field theory trades microscopic details for self-consistency. Instead of tracking every interaction, it asks whether the average activity assumed as input equals the average activity produced as output. Static mean field answers this for stationary means. Dynamic mean field asks a harder question: whether the time-dependent fluctuations or autocorrelations assumed as input are reproduced by the effective representative unit or population.

\paragraph{Main mechanism: deriving the static mean-field equation.}
Consider the Ising energy
\[
E(s)=-\frac12\sum_{i,j}J_{ij}s_i s_j-\sum_i h_i s_i.
\]
In mean field, approximate the influence of other variables on \(S_i\) by replacing \(S_j\) with \(m_j\). The effective energy contribution for \(s_i\) is
\[
E_i(s_i)=-s_i\left(h_i+\sum_j J_{ij}m_j\right).
\]
Let
\[
H_i^{\mathrm{mf}}=h_i+\sum_j J_{ij}m_j.
\]
Then the approximate distribution over \(s_i\) is
\[
q_i(s_i)=\frac{\exp(\beta s_iH_i^{\mathrm{mf}})}
{2\cosh(\beta H_i^{\mathrm{mf}})}.
\]
Its mean is
\[
m_i=\mathbb E_{q_i}[S_i]
=\frac{e^{\beta H_i^{\mathrm{mf}}}-e^{-\beta H_i^{\mathrm{mf}}}}
{e^{\beta H_i^{\mathrm{mf}}}+e^{-\beta H_i^{\mathrm{mf}}}}
=\tanh(\beta H_i^{\mathrm{mf}}).
\]
Thus
\[
m_i=\tanh\left(\beta\left[h_i+\sum_jJ_{ij}m_j\right]\right).
\]
The approximation is the replacement of interactions by a self-consistent average field.

\paragraph{Worked examples.}
\emph{Small numerical example.}
For one mean-field unit with \(h=0.2\), \(J=1\), \(\beta=1\), the self-consistency equation is
\[
m=\tanh(0.2+m).
\]
Starting from \(m_0=0\), one fixed-point iteration gives
\[
m_1=\tanh(0.2)\approx0.197,
\]
\[
m_2=\tanh(0.397)\approx0.378.
\]
The iteration illustrates self-consistency: the mean creates a field, and the field creates a new mean.

\emph{Computational neuroscience example.}
A ring-attractor neural field for head direction can use \(x\) as preferred direction and a kernel \(w(x-x')\) with local excitation. A localized bump of \(u(x,t)\) represents the current head direction. Perturbing the bump and watching it return, move, or spread tests the model's stability.

\emph{AI example.}
Variational inference often uses a mean-field family
\[
q(z)=\prod_i q_i(z_i)
\]
to approximate a hard posterior. Energy-based models and Boltzmann machines face hard partition functions, so mean-field, Monte Carlo, or contrastive approximations are used. Static mean field estimates average activations; dynamic mean-field ideas appear when random recurrent networks are approximated by effective stochastic processes.

\paragraph{Common misconceptions and failure modes.}
Mean field is not the same as taking an ordinary sample mean. It is a self-consistent approximation to interacting variables. It may fail when correlations, clusters, rare states, or strong spatial structure are essential. A neural field is continuous in space by modeling choice; it does not mean cortex is literally continuous at every scale. Dynamic mean field is not just a time-dependent static mean; it must close statistics of fluctuations, not only average rates.

\paragraph{Exercises.}
\begin{enumerate}[label=\textbf{49.4.\arabic*.}, leftmargin=*]
\item \textbf{Mechanical.} Identify the leak, recurrent integral, nonlinearity, and input terms in the neural field equation.
\item \textbf{Proof or derivation.} Derive \(m_i=\tanh(\beta[h_i+\sum_jJ_{ij}m_j])\) from the one-site mean-field distribution \(q_i(s_i)\).
\item \textbf{Numerical/computational.} Perform three fixed-point iterations for \(m=\tanh(0.5m)\) starting from \(m_0=0.8\).
\item \textbf{Interpretation.} Explain the difference between static mean field and dynamic mean field in a recurrent neural population.
\item \textbf{Debug the misconception.} A student says, "Mean-field theory proves correlations do not matter." Correct the statement by explaining that mean field approximates or neglects correlations and must be checked.
\end{enumerate}

\paragraph{Closing bridge.}
Recurrent circuits give mechanisms for stable activity, memory, oscillation, and spatial organization. Energy and free-energy ideas connect these circuits to statistical physics and modern energy-based AI, while mean-field approximations explain why average population descriptions can sometimes replace microscopic interactions. Chapter 50 shifts from circuit mechanisms to inference and action, asking how Bayesian updates, evidence accumulation, predictive coding, and active inference can be stated carefully as mathematical models of perception and decision-making.

\clearpage
